
En traitement du signal, la phase d'un signal est intrinsèquement lié à la notion de fréquence instantanée, qui joue un rôle important en analyse temps-fréquence. 
C'est donc de point que commencera notre discussion pour introduire la phase géométrique.
Pour cela, seront rapidement introduit quelques notions et résultats d'analyse temps-fréquence dans le cas univarié (sec. \ref{subsec:ana_temp-freq}). Suite à quoi sera définie la phase instantanée pour le cas multivarié (sec. \ref{subsec:intro_phased}), qui permettre enfin de mettre en évidence la phase géométrique (sec. \ref{subsec:intro_phaseg}).
\\

\thoughts{OUTDATED} Dans une seconde partie, seront introduit les signaux bivarié dit AM-FM-PM  (sec. \ref{subsec:AM-FM-PM}), dont la phase géométrique sera calculée explicitement (sec. \ref{subsec:phase_g2AM-FM-PM}), ce qui permettra de mettre en évidence certaines de ses propriétés. Dans une dernière section, seront proposées des généralisations des signaux AM-FM-PM au delà du cas bivarié et seront discutées leur pertinence quant à l'étudie de la phase géométrique (sec. \ref{subsec:gene_AM-FM-PM}). 

\section{Introduction de la phase géométrique} \label{sec:intro_phaseg}




\subsection{Cas univarié : signaux AM-FM} \label{subsec:ana_temp-freq}

%\subsubsection{Quelque notion d'analyse temps-fréquence}

En traitement du signal, l'analyse fréquentielle par la transformée de Fourier est un incontournable. 
Seulement, cette transformation fait perdre toute notion temporelle : si l'étude du spectre du signal permet de dire quelles fréquences apparaissent dans le signal, elle ne permet pas de dire à quel(s) moment(s). 
C'est en réponse à cela, entre autre, qu'est développé l'analyse temps-fréquence. 
\\
À cette fin, sont définies les paramètres instantanées d'un signal :\par

\begin{definition}[Paramètres instantanées] \label{def:param_instant}
	Soit $x$, est un signal complexe écrit sous forme exponentielle :
	\begin{align}
		x\ &:\ \begin{aligned}\R\ &\lr\qquad \C \\
			t\ &\longmapsto\ a(t)e^{i\phi(t)}
		\end{aligned}  &  \text{où }\quad a(t)\in\R^+\quad &\text{et}\quad \phi(t)\in\R
	\end{align}
	$a$ est appelé \emph{amplitude instantanée} du signal, $\nicefrac{1}{2\pi}\phi'$ sa \emph{fréquence instantanée} et sa \emph{phase instantanée} est définie --- modulo un choix de phase initiale --- par :
	\begin{equation} \label{eq:phasei}
		\phasei(x,t_0,t) = \phi(t) - \phi(t_0)
	\end{equation}
\end{definition}
\skipl 

Pour les signaux réels, ces notions sont moins évidentes à définir puisqu'elles demandent d'écrire les signaux sous la forme :
\[x(t) = a(t) \cos\phi(t)\]
\\
Auquel cas, le choix de la pair $(a,\phi)$ n'est pas unique. Il existe tout de même un ``bon'' choix de telle pair dans le cas des signaux dit AM-FM :
\begin{definition}[Signal AM-FM]
	Un signal réel de la forme :
	\begin{align}
		x\ &:\ \begin{aligned}\R\ &\lr\qquad \R \\
			t\ &\longmapsto\ a(t) \cos\phi(t)
		\end{aligned}  &  \text{où }\quad a(t)\in\R^+\quad
	\end{align}
	est dit \emph{AM-FM} (\emph{amplitude and frequency modulated}) si $a$ et $\cos\phi$ admettent des transformée de Fourier et si de plus la première a un spectre concentré sur les bases fréquences, la seconde concentré sur les hautes fréquences et que les deux ne se chevauche pas.
	Formellement, ces conditions demande qu'il existe $\lambda\in\R^+$ telle que :
	\begin{equation}\label{eq:condi_AM-FM}
		\supp \Fou{a} \subset [-\lambda, \lambda],\quad \supp \Fou{\cos\phi} \subset \R\setminus[-\lambda,\lambda]
	\end{equation}
	\\
	Dans ce cas, $a$ et $\phi$ donne lieu au même vocabulaire que pour le cas complexe (\cref{def:param_instant}).
\end{definition}
\skipl
\\
Ces conditions sont liées au théorème de Bedrosian, et plus de détail se trouve dans l'annexe \ref{ann:complement_t-f}. Pour le dire rapidement, elles exigent que toutes les hautes fréquences de $x$ se trouvent dans la phase. Intuitivement, cela évite que toute les fréquences aillent dans l'amplitude $a$, auquel cas, $x$ n'aurait ``pas de fréquence'' au sens où $\phi$ pourrait être choisie constante, voir nulle.
\\
Sous ces conditions, $x$ peut être vu comme le signal complexe $\SA{x}$ telle que :
\[\SA{x}(t) = a(t) e^{i\phi(t)}= a(t)\cos\phi(t) + ia(t)\sin\phi(t) = x  + i\Im m \SA{x}\]
L'on parle alors de transformée en \emph{signal analytique} de $x$ et $\SA{x}$ a, par construction, les mêmes paramètres instantanée que $x$.
Là encore, le lecteur est renvoyé vers l'annexe \ref{ann:complement_t-f} pour plus de détail.
\\


%\subsubsection{Le problème des signaux multivariés}

L'intérêt d'introduire toutes ces notions est que les signaux multivariés --- même complexe --- souffre du même problème que les signaux réels. 
En effet, en écrivant un signal $\x$ sous la forme :
\[\forall t\in\R,\qquad 
\x(t) = \begin{pmatrix} A_1(t)e^{i\phi_1(t)} \\ A_2(t)e^{i\phi_2(t)} \\ \vdots \\ A_n(t)e^{i\phi_n(t)}
\end{pmatrix}\]
\\
le fait que $\x$ soit à valeur dans $\C^n$ impose un choix naturel d'amplitude instantanée : sa norme. Pour ce qui est de la phase instantanée, en revanche, n'importe qu'elle choix de $\phi$ convient \apriori~:
\[\forall t\in\R,\qquad 
\x(t) = \begin{pmatrix} A_1(t)e^{i\phi_1(t)} \\ A_2(t)e^{i\phi_2(t)} \\ \vdots \\ A_n(t)e^{i\phi_n(t)} \end{pmatrix}
= a(t)e^{i\phi(t)}\begin{pmatrix} a_1(t)e^{i\psi_1(t)} \\ a_2(t)e^{i\psi_2(t)} \\ \vdots \\ a_n(t)e^{i\psi_n(t)} \end{pmatrix}
\qquad\text{ avec }\qquad 
\left\{ \begin{aligned}
	& a(t) = \| \x(t) \|_2 \\
	& \big\|(a_i)_{1\leq i\leq n}\big\|_2 = 1 \\
	& \phi_i = \phi + \psi_i \end{aligned}\right.\]
\\
Il suffit que les $\psi_i$ soient ajustés pour assurer que $\ \phi_i = \phi + \psi_i$.
\\
\begin{remarque}
	À noter, que si $a$ et $\phi$ sont correspondent respectivement à une amplitude et une phase, le vecteur restant $\big( a_ie^{\phi_i} \big)_{1\leq i\leq n}$ correspond à un vecteur de polarisation, sur lequel nous reviendrons dans la \cref{sec:AM-FM-PM} suivante.
\end{remarque}



\subsection{Phase et fréquence instantanée de signal multivarié }\label{subsec:intro_phased}

On se propose ici de définir la phase instantanée comme suit :
\begin{definition}[Phase dynalique/instantanée] \label{def:phase_d}
	La \emph{phase instantanée} ou \emph{dynamique} (à l'instant $t$ partant de $t_0$) d'un signal multivarié $\x = a\big(a_ie^{i\phi_i}\big)_{1\leq i\leq n} \in \conti[1]{\R}{\C^n}$ quelconque, est définie par la formule :
	\begin{equation} \label{eq:phase_d}
		\forall t_0, t\in\R, \quad \phased(\x, t_0,t) \defeq \int_{t_0}^t \frac{\Im m \big\langle \dot{\x}(s) , \x(s) \big\rangle}{\|\x(s)\|^2} ds =\sum_{i=1}^n \int_{t_0}^t a_i(s)^2 \phi_i'(s)ds
	\end{equation}
	On s'autorisera à omettre les paramètres de $\phased$ lorsque cela ne prête pas à confusion.
\end{definition}

\begin{remarque}
	Outre l'aspect variationnelle de cette formule, le terme ``dynamique'' viens du fait que, lorsque $\x$ suit une équation de Schrödinger :
	\begin{equation}\label{eq:schrodinger}
		i\frac{d \x(t)}{dt} = h\x(t)
	\end{equation}
	la dérivée $\dot{\x}$ dans la formule \eqref{eq:phase_d} ci-dessus se voit remplacé par l'hamiltonien $h\x$ {\normalfont \cite[sec. 2]{bohm_geometric_2003}, \cite[p.~215]{mukunda_quantum_1993}}, donnant :
	\[\phased' = -i \int_{t_0}^t\frac{\big\langle h\x(s) , \x(s) \big\rangle}{\|\x(s)\|^2} ds\] 
	\\
	Sachant que $\x$ n'a aucune raison de suivre une telle équation dans notre cas, poser $h = i\frac{d}{dt}$ enlève toute contrainte, auquel cas $\phased'=\phased$.
\end{remarque}
\skipl

Ceci étant, deux arguments sont donnés pour motiver cette définition :
\\

\subsubsection*{Argument variationnelles}

Le premier, fortement inspiré par les travaux de Lilly \& Olhede  \cite{lilly_analysis_2012}, consiste à généraliser la condition \eqref{eq:condi_AM-FM} de séparation haute/basse fréquences sur les signaux AM-FM.
Pour cela, l'on commence par faire apparaître une phase $\phi$ --- pour l'instant inconnue --- en écrivant $\x$ sous la forme :
\[\forall t\in\R,\qquad \x(t) = e^{i\phi(t)} e^{-i\phi(t)} \x(t) \defeq e^{i\phi(t)} \bf{y}(t)\]
\\
Si $\phi$ est bien choisie, alors $\bf{y}$ ne devrait contenir que les informations associées à l'amplitude et la polarisation de $\x$. Or, conformément à la condition \eqref{eq:condi_AM-FM}, la phase doit contenir les hautes fréquences du signal et, inversement, les basses fréquences doivent se trouver dans le reste. 
\\
La fréquence donnant, pour le dire vite, la vitesse d'ondulation, la contrainte sur $\x$ va être de limite les variations de  $\bf{y}$. Concrètement, $\phi$ doit être choisie de sorte à minimiser la dérivée $\dot{\bf{y}}$ :
\[\forall t\in\R,\qquad \phi(t) = \argmin{\theta(t)}{\big\|\dot{\bf{y}}(t)\big\|_2}^2 = \argmin{\theta(t)}{\Big\|e^{-i\theta(t)}\big(\dot{\x}(t) - i\theta'(t) \x(t)\big) \Big\|_2}^2 = \argmin{\theta(t)}{\big\|\dot{\x}(t) - i\theta'(t)\x(t)\big\|_2}^2\]
\\
La contrainte ne dépendant que de la dérivée $\theta'$, on se ramène à :
\[\min_{\theta(t)}{\|\dot{\bf{y}}(t)\|_2}^2 = \min_{\theta'(t)}{\big\|\dot{\x}(t) - \theta'(t) \x(t)\big\|_2}^2\]
\\
En rappelant que $\frac{d}{dx}{\big\|f(x)\big\|_2}^2 = 2\Re e\big\langle f(x), f'(x)\big\rangle$, il vient que ce minimum\footnote{\itshape
	L'extremum obtenu est l'unique minimum globale puisque $t\longmapsto \|at + b\|^2$ est strictement convexe pour $a\neq0$.}
est atteint par $\phi'(t)$ à condition que :
\begin{align*}
	\frac{d}{d\phi'}{\big\| \dot{\x} - i\phi' \x\big\|_2}^2 = 0 \quad \Llr\quad
	0 &= 2\Re e\left\langle  \dot{\x} - i\phi' \x ,  \frac{d}{d\phi'}\big(\dot{\x} - i\phi' \x\big)\right\rangle \\
	&= 2\Re e\big\langle  \dot{\x} - i\phi' \x ,  - i \x\big\rangle \\
	&= 2\Re e\Big(i\big\langle  \dot{\x} ,  \x\big\rangle\Big) + 2\phi'\Re e\big\langle   \x ,  \x\big\rangle\\
	&= -2\Im m\big\langle  \dot{\x} ,  \x\big\rangle + 2\phi'{\| \x\|_2}^2
\end{align*}
Ainsi $\displaystyle \ \phi' = \frac{\Im m\big\langle  \dot{\x} ,  \x\big\rangle}{{\| \x\|_2}^2}\ $ et :
\begin{equation}\label{eq:phas_inst_v1}
  \phi(t) = \Im m\int_{t_0}^t \frac{\big\langle \dot{\x}(s) , \x(s) \big\rangle}{\|\x(s)\|^2} ds = \phased(\x,t_0,t)
\end{equation}
\\

\subsubsection*{Arguments des moyennes}

Autre argument, cette fois inspiré de \cite{cano_mathematical_2022}, ce base sur la notion de fréquence moyenne.
D'abord dans le cas d'un signal complexe univarié, sont définies les fonctions de densités d'énergie (resp. d'énergie spectale) comme :
\begin{align}\label{eq:densi_dE}
	\densit\ &:\quad \begin{aligned}\R\ &\lr\quad \R^+ \\ t\ &\longmapsto\ \big|x(t)\big|^2 \end{aligned}  
	&
	\text{resp.}\qquad \densis\ &:\quad \begin{aligned}\R\ &\lr\quad \R^+ \\ \nu\ &\longmapsto\ \big|\fou{x}(\nu)\big|^2 \end{aligned}
\end{align}
\\
À partir de ces dernières est définie la fréquence moyenne de $x$ comme comme l'espérance $\esp[\densis]{\nu}$ de $\densis$. Cette fréquence moyenne est lié à la fréquence instantanée par la formule :\footnote{cette formule de généralise à tout les moments de $\densis$ et existe également pour les moments de $\densit$, voir \cite[sec. 1.4]{cohen_time_1995} pour une démonstration ``à la physicienne'' \thoughts{... ou bien en annexe ?}}
\begin{equation}\label{eq:esp_freq}
	\esp[\densis]{\nu} = \frac{1}{2\pi}\int_\R \phi'(t)\densit(t)dt = \frac{1}{2\pi} \esp[\densit]{\phi'}
\end{equation}
\\
Dans le cas d'un signal $\x=(x_i)_{1\leq i\leq n}$ multivarié, les densités d'énergies se définissent comme :
\begin{align*}%\label{eq:densi_dEi}
	\densit_i\ &:\quad \begin{aligned}\R\ &\lr\quad \R^+ \\ t\ &\longmapsto\ \big|x_i(t)\big|^2 = a(t)^2 a_i(t)^2 \end{aligned}  
	&
	\densis_i\ &:\quad \begin{aligned}\R\ &\lr\quad \R^+ \\ \nu\ &\longmapsto\ \big|\fou{x}_i(\nu)\big|^2 \end{aligned} \\ \\
	%\label{eq:densi_dE-mv}
	\densit\ &:\quad \begin{aligned}\R\ &\lr\quad \R^+ \\ t\ &\longmapsto\ \big\|\x(t)\big\|^2 = \sum_{i=1}^n \densit_i(t) \end{aligned}  
	&
	\densis\ &:\quad \begin{aligned}\R\ &\lr\quad \R^+ \\ \nu\ &\longmapsto\ \big\|\fou{\x}(\nu)\big\|^2 = \sum_{i=1}^n \densis_i(t) \end{aligned}	
\end{align*}
Le second argument consiste alors à dire que l'égalité des moments $\eqref{eq:esp_freq}$ doit rester vraie dans le cas multivarié. Cela assure au moins que la fréquence instantanée de $\x$, $\nicefrac{1}{2\pi}\phi'$, à pour moyenne $\esp[\densis]{\nu}$, la fréquence moyenne en sens de Fourier.
\\

En appliquant la formule \eqref{eq:esp_freq} au $\densis_i$, et en notant toujours $\x = a\big(a_ie^{i\phi_i}\big)_{1\leq i\leq n}$, on obtient :
\begin{align*}
	\esp[\densis]{\nu} = \int_\R \nu\densis(\nu)d\nu &= \int_\R \nu\sum_{i=1}^n \densis_i(\nu) d\nu \\
	&= \sum_{i=1}^n\esp[\densis_i]{\nu} \\
	&= \sum_{i=1}^n\frac{1}{2\pi}\int_\R \phi_i'(t)\densit_i(t)dt \\
	&= \frac{1}{2\pi}\int_\R a(t)^2\sum_{i=1}^n\phi_i'(t)a_i(t)^2 dt 
	\\ &= \frac{1}{2\pi} \esp[\densit]{\sum_{i=1}^n \phi_i'{a_i}^2}
\end{align*}
\\
Ce qui mène à poser $\displaystyle \ \sum_{i=1}^n \phi_i'(t){a_i}^2(t)\ $ pour la fréquence instantanée, avec la phase associée :
\begin{equation}\label{eq:phas_inst_v1}
	\phi = \int_{t_0}^t \sum_{i=1}^n \phi_i'(s){a_i}(s)^2ds 
	= \sum_{i=1}^n \int_{t_0}^t \phi_i'(s){a_i}(s)^2ds 
	%= \sum_{i=1}^n \esp[\nicefrac{\densit_i}{\densit}]{\phi_i'}
\end{equation}
\\

Formule qui concorde bien avec celle de la phase dynamique une fois explicité :
\begin{align*}
	\Im m\frac{\big\langle \dot{\x}(t) , \x(t) \big\rangle}{\|\x(t)\|^2} &= \Im m\left( \frac{1}{a(t)^2} \sum_{i=1}^n \Big( \big(aa_i\big)'(t) +a(t)a_i(t)i\phi_i'(t)\Big)e^{i\phi_i(t)}\congu{a(t)a_i(t)e^{i\phi_i(t)}} \right) \\
	&=\frac{1}{a(t)^2}  \Im m\left( \sum_{i=1}^n a(t)a_i(t)\big(aa_i\big)'(t) +ia(t)^2a_i(t)^2\phi_i'(t) \right) \\
	&= \frac{1}{a(t)^2} \sum_{i=1}^n a(t)^2a_i(t)^2 \phi_i'(t) \\
	&= \sum_{i=1}^n a_i(t)^2 \phi_i'(t)
\end{align*}
D'où
\[\Im m\int_{t_0}^t \frac{\big\langle \dot{\x}(s) , \x(s) \big\rangle}{\|\x(s)\|^2} ds = \int_{t_0}^t \sum_{i=1}^n a_i(s)^2 \phi_i'(s) = \sum_{i=1}^n \int_{t_0}^t a_i(s)^2 \phi_i'(s)ds\]
\\



\subsection{Apparition de la phase géométrique}\label{subsec:intro_phaseg}

Cela étant dit, il existe une autre façon, plus simple, d'obtenir la phase d'un signal. D'abord, dans le cas univarié, la phase instantanée de $x=ae^{i\phi}$ peut être réécrite comme :
\[\phi(t)-\phi(t_0)  = \arg\left( x(t) \congu{x(t_0)} \right)\]
\\
Formule qui se généralise en cas multivarié par ce qui sera appelé la \emph{phase totale} du signal :
\begin{equation}\label{eq:phase_t}
	\phaset(\x, t_0, t) \defeq \arg\big\langle \x(t), \x(t_0)\big\rangle
\end{equation}
\\
D'un point de vu géométrique, il est bien connue que le produit scalaire entre deux vecteurs réels $u,v\in\R^n$ est lié à l'angle $\angle(u,v)$ entre ces derniers par la formule :
\[\langle u,v\rangle_\R = \|u\|^2 \|v\|^2 \cos \angle(v,u)\]
\\
Pour le produit hermitien, cet angle ce retrouve dans l'argument, de sorte que si $u$ et $v$ sont complexes :
\[\langle u,v\rangle_\C = \|u\|^2 \|v\|^2 e^{i \angle(v,u)}\]
\\
En ce sens, la phase totale calcul explicitement l'angle entre $\x(t_0)$ et $\x(t)$. La question est alors de savoir si $\phased$ correspond à cette angle. 
Un calcul explicite montre que c'est bien le cas en univarié ; en notant $\x = ae^{i\phi}$, il vient :
\begin{align*}
	\phased(\x) = \Im m\int_{t_0}^t \frac{\big\langle \dot{\x}(s) , \x(s) \big\rangle}{\|\x(s)\|^2} ds &= \Im m \int_{t_0}^t \frac{\big(a'(s) + ia(s)\phi'(s) \big) e^{i\phi(s)} \congu{a(s) e^{i\phi(s)}}}{a^2(s)} ds \\
	&= \int_{t_0}^t \frac{a^2(s)\phi'(s))}{a^2(s)} ds \\
	&= \phi(t) - \phi(t_0)
\end{align*}
\skipl

Dans le cas multivarié, en revanche, c'est une autre histoire. En notant cette fois le signal $\ \x = ae^{i\phased} \big( a_ie^{\psi_i} \big)_{1\leq i\leq n}$, la phase totale se réécrit :
\begin{equation}\label{eq:diff_phases_t/d}
	\begin{aligned}
		\phaset(\x,t_0, t) &= \arg \left(a(t)a(t_0) e^{i\big(\phased(t) - \phased(t_0)\big)}\sum_{i=1}^n a_i(t)a_i(t_0)e^{i(\psi_i(t)-\psi_i(t_0))} \right) \\
		&= \phased(t) + \arg \left(\sum_{i=1}^n a_i(t)a_i(t_0)e^{i(\psi_i(t)-\psi_i(t_0))} \right)  \qquad\qquad\qquad\qquad \text{car } \phased(t_0,t_0) = 0
		%\\ &= \phased + \arctan \left( \frac{\sum_i a_i(t)a_i(t_0)\sin\big( \psi_i(t)-\psi_i(t_0)\big)}{\sum_i a_i(t)a_i(t_0)\cos\big( \psi_i(t)-\psi_i(t_0)\big)}  \right)
	\end{aligned}
\end{equation}
\\
Apparaît alors un terme de déviation de la phase dynamique par rapport à la phase totale, qui (surprise) est appelé phase géomatique et sera notée :
\begin{equation}\label{eq:phase_g}
	\phaseg(\x,t_0,t) \defeq \phaset(\x, t_0,t) - \phased(\x, t_0,t)
\end{equation}
Déviation qui s'observe expérimentalement, comme le montre la \Cref{fig:calc_diff_phases} ci-dessous.
\\
\begin{figure}[h]
	\includegraphics[width=0.6\textwidth]{fig/placeholder}
	\caption[Déviation de la phase dynamique par rapport à la phase totale]{Sur le graphe de gauche, le signal $\x$ à valeur dans $\R^2$ et dans celui de droite le calcul des phases dynamique et totale ainsi que de leur différence. Résultat tiré des simulation de Le Bihan \etal~\cite{le_bihan_modephysiques_2023}}
	\label{fig:calc_diff_phases}
\end{figure}
\\

Un résultat bien connue en physique \cite{bohm_geometric_2003,mukunda_quantum_1993,nakahara_geometry_2003} est que cette troisième phase est invariante par transformation de jauge. Dans notre contexte, cela signifie que si $\x$ et $\Tilde{\x}$ sont deux signaux multivarié complexe tel que $\ \Tilde{\x} = e^{i\alpha}\x$, avec $\alpha$ une \underline{fonction} continue du temps, alors :
\begin{align*}
	\phaseg(\Tilde{\x}) &= \phaset(\Tilde{\x}) - \phased(\Tilde{\x})  = \phaset(\x) - \phased(\x) = \phaseg(\Tilde{\x})\\
	%&= \arg\big\langle \Tilde{\x}(t), \Tilde{\x}(t_0)\big\rangle - \Im m\int_{t_0}^t \frac{\big\langle \dot{\Tilde{\x}}(s) , \Tilde{\x}(s) \big\rangle}{\|\Tilde{\x}(s)\|^2} ds \\
	%&= \arg\big\langle e^{i\alpha(t)}\x(t), e^{i\alpha(t_0)}\x(t_0)\big\rangle - \Im m\int_{t_0}^t \frac{\big\langle \dot{e^{i\alpha(s)}\x}(s) , e^{i\alpha(s)}\x(s) \big\rangle}{\|e^{i\alpha(s)}\x(s)\|^2} ds 
\end{align*}
\\
Cela signifie plusieurs choses. 
D'une part, le fait cette phase soit invariante pas transformation de jauge montre qu'elle est associée à la composante de $\x$ dont il a été mentionné (\textit{sec.} \ref{subsec:ana_temp-freq}) qu'elle était associée à sa polarisation. 
Ce point sera discuté dans la \cref{sec:AM-FM-PM} suivante.
 \\
D'autre part, cela suggère que la phase géométrique doit avoir une certaine robustesse au bruit. \thoughts{Pas sur de ça :} Pour argumenter ce point, rappelons que le bruit d'un signal se trouve généralement dans les hautes fréquences de ce dernier.
Or, en admettant, comme ça l'a été abordé dans la \cref{subsec:intro_phased} précédent, que la phase dynamique contient les hautes fréquences, il y a d'autant plus de chances que cette phase est un bon comportement par rapport au bruit.
\\
Ainsi, s'il est possible de tirer des informations de $\phaseg$, alors ces informations doivent pouvoir être obtenue même pour des signaux particulièrement bruités, comme c'est le cas des mesures d'ondes gravitationnelles \thoughts{(j'aimerais beaucoup étudier ça mais ce sera sûrement pour une thèse, si thèse il y a)}.
\\

Cela étant dit, avec le calcul \eqref{eq:diff_phases_t/d} précédent, il peut sembler que le travail sur la phase géométrique est terminée en cela qu'une formule explicite est donnée. Deux remarques à ce sujet :
\\
D'abord, cette formule demande de connaître les $\psi_i$, qui eux-mêmes sont obtenue en extrayant la phase dynamique au signal. Or, la formule de $\phased$ n'est pas la plus appropriée au traitement du signal puisque qu'elle fait intervenir intégral et dérivée.
\\
Aussi mais surtout, cette formule cache la profondeur dernière l'étude de la phase géométrique, à commencer par l'origine de son nom.
\\
\thoughts{À modifier en fonction de la suite :}
Dans ce mémoire le calcul pratique de la phase géométrique ne sera pas aborder pour, à la place, se pencher sur le second point : l'aspect géométrique d'une phase éponyme.
\\ \\

\thoughts{Est-ce que je dois parler de l'espace projectif ou est-ce que je le garde pour la grande partie II ?}
\\



\section{Première étude : cas des signaux AM-FM-PM bivariés} \label{sec:AM-FM-PM}

Pour une première étude de la phase géométrique du signal, Le Bihan \etal~se sont penchés sur un cas particulier de signal bivarié \cite{flamant_timefrequency_2019,le_bihan_modephysiques_2023, le_bihan_geometric_2024}. Ces signaux, dit AM-FM-PM, sont présenté dans une première partie et le calcul explicite de leur phases --- totale, dynamique et géométrique --- est présenté. Dans une seconde partie... \thoughts{je sais pas trop}
\\



\subsection{Définitions et calcul des phases} \label{subsec:AM-FM-PM}

Ces signaux AM-FM-PM viennent généraliser les signaux AM-FM univarié en tenant compte de l'état de polarisation permis par le passage au 2D. 
En quelques mots, dans le cas le plus simple, un signal bivarié à valeurs réelles $s$ va décrire une ellipse en cours du temps. 
On parle de polarisation elliptique et $s$ va s'écrire :
\[s(t) = a \begin{pmatrix} \cos\theta & -\sin\theta \\ \sin\theta  &  \cos\theta \end{pmatrix} \begin{pmatrix} \cos\chi \cos\varphi(t) \\ \sin\chi \sin\varphi(t) \end{pmatrix}  \qquad \text{ où }\quad  a\in\R^+,\ \theta \in \left]-\frac{\pi}{2}, \frac{\pi}{2}\right],\ \chi \in \left[-\frac{\pi}{4}, \frac{\pi}{4}\right] \]
\\
Les paramètres $a$ et $\chi$ caractérisent respectivement la taille et l'excentricité de l'ellipse, $\theta$ son orientation dans le plan et $\varphi(t)$ précise où se trouve $s$ à l'instant $t$ sur cette ellipse. 
Le tout est représenté sur la \Cref{fig:ellipse1polar} ci-dessous :
\begin{figure}[h]
	\includegraphics[width=0.45\textwidth]{fig/placeholder}
	\caption[Ellipse de polarisation d'un signal bivarié réel]{Ellipse de polarisation du signal $s$ sur laquelle sont représenter ses paramètres $a,\varphi,\theta,\chi$.}
	\label{fig:ellipse2polat}
\end{figure}
\\
En autorisant les paramètres de polarisation à varier au cours du temps et après une transformation en signal analytique, mentionnée dans la \cref{subsec:ana_temp-freq}, on obtient la définition suivante :
\\
\begin{definition}[Signal AM-FM-PM] \label{def:AM-FM-PM}
	Un signal bivarié complexe $\x$ \emph{AM-FM-PM} (\emph{amplitude, frequency and polarization modulated}) est caractérisé par quatre paramètres $a,\varphi,\theta$ et $\chi$, respectivement à valeur dans $\R^+$, $\R$, $]-\frac{\pi}{2}, \frac{\pi}{2}]$ et $[-\frac{\pi}{4}, \frac{\pi}{4}]$, vérifiant :
	\begin{align}\label{eq:condi_AM-FM-PM}
		\big| \varphi'(t) \big| &\gg \big| \theta'(t) \big| ,\ \big| \chi'(t) \big| ,\ \left| \frac{a'(t)}{a(t)}\right|  &  \left| \frac{\varphi'(t)}{\varphi(t)}\right| \gg 1
	\end{align}
	Auquel cas, $\x$ prend la forme, $\forall t\in\R$ :
	\begin{equation}\label{eq:AM-FM-PM}
		\x(t) = a(t)e^{i\varphi(t)} R_{\theta(t)} \begin{pmatrix} \cos\chi(t) \\ -i\sin\chi(t) \end{pmatrix} 
		= a(t)e^{i\varphi(t)} \begin{pmatrix} \cos\theta(t) \cos\chi(t) + i\sin\theta(t) \sin\chi(t) \\ \sin\theta(t) \cos\chi(t) - i\cos\theta(t) \sin\chi(t) \end{pmatrix}
	\end{equation}
	où $R_{\theta(t)}$ est la matrice de rotation d'angle $\theta(t)$.
\end{definition}
\skipl

La transformée en SA permet de transformer les signaux à valeurs réelles en complexes, ce qui est nécessaire\footnote{\itshape
	Nous reviendrons sur ce point dans la dernière partie de mémoire \thoughts{... est-ce qu'on le fait vraiment?}}
pour étudier la phase géométrique car c'est uniquement dans le cadre de complexe qu'elle a été étudiée jusqu'à présent. Et, comme pour les signaux AM-FM, les hypothèses sur $a,\varphi,\theta,\chi$ assure que les paramètres soient interprétables comme sur la \Cref{fig:ellipse2polat} précédente.
\\

Les trois phases de tels signaux sont alors données par la \cref{prop:phases_2var} suivante :
\begin{proposition}[phases de signal AM--FM--PM]\label{prop:phases_2var}
	Les trois phases d'un signal bivarié AM--FM--PM $\x$ de paramètres $(a,\varphi,\theta,\chi)$ sont données par les formules :
	\begin{equation}\label{eq:phased_2var}
		\phased(\x, t_0,t) = \varphi(t) -\varphi(t_0) + \int_{t_0}^t\theta'(s) \sin2\chi(s) ds
	\end{equation}
	\begin{equation}\label{eq:phaset_2var}
	\begin{aligned}
		\phaset(\x,t_0,t) &= \varphi(t)-\varphi(t_0) + \arg\Big(\cos\Delta\theta \cos\Delta\chi + i\sin\Delta\theta \sin\big(\chi(t_0)+\chi(t)\big)\Big) \\
		&= \varphi(t)-\varphi(t_0) + \arctan\left(\tan\Delta\theta \frac{ \tan\chi(t_0)+\tan\chi(t)}{1 + \tan\chi(t_0)\tan\chi(t)}\right)
	\end{aligned}
	\end{equation}
	\begin{equation}\label{eq:phaseg_2var}
	\begin{aligned}
		\phaseg(\x,t_0,t) &= \phaset(\x,t_0,t) - \phased(\x,t_0,t) \\
			&= \arctan\left(\tan\Delta\theta \frac{ \tan\chi(t_0)+\tan\chi(t)}{1 + \tan\chi(t_0)\tan\chi(t)}\right) - \int_{t_0}^t\theta'(s) \sin2\chi(s) ds
	\end{aligned}
	\end{equation}
	\\
	où $\ \Delta y = y(t) - y(t_0)\ $ pour $\ y = \chi$. La démonstration se trouve en annexe \ref{ann:demo_phases_2var}
\end{proposition}
\skipl

Deux remarques sur ces formules. 
La première est que la phase géométrique ne dépend que des paramètres de polarisations $\theta$ et $\chi$, un résultat rassurant sachant qu'elle est invariantes pas transformation de jauge.
La seconde, nettement plus troublante, est que $\varphi$ ne s'interprète ni comme la phase totale ni comme la phase dynamique. Plus loin, \thoughts{\cref{subsec:aller_plus_loin}}, nous reviendrons sur laquelle des deux ``doit'' représenter $\varphi$.
\\

Aussi, la phase géométrique est séparée en deux termes. Le premier, en arctangente, correspond à l'angle formé par les vecteurs portant la polarisation instantanée, c'est-à-dire les vecteurs :
\begin{align*}
	R_{\theta(t_0)} \begin{pmatrix} \cos\chi(t_0) \\ -i\sin\chi(t_0) \end{pmatrix}\qquad \text{ et }\qquad R_{\theta(t)} \begin{pmatrix} \cos\chi(t) \\ -i\sin\chi(t) \end{pmatrix}
\end{align*}
\\

Pour interpréter le second membre, c'est beaucoup moins clair. Il faut d'intéresser  la projection  $\ \rho_{\x} = \frac{1}{\|\x \|^2} \congu{\x }\, ^t\x $. Avec des calculs (interminables), on obtient l'écriture :
\[\rho_{\x} = id + \sin2\theta\cos2\chi \sigma_1 + \sin 2\chi \sigma_3 + \cos2\theta \cos2\chi \sigma_3\]


\begin{wrapfigure}{r}{0.3\textwidth}
	\begin{tikzpicture}[scale=0.8]
	%\draw[black, opacity=0.1] (-4.5,-4.5) grid (4.5,4.5);
	\draw[-Stealth] (-2.5,0) -- (4.5,0) node[below]{$\sigma_1$};
	\draw[-Stealth] (0,-4.5) -- (0,4.5) node[right]{$\sigma_2$};
	\draw[-Stealth] (2.5,1.2) -- (-2.5,-1.2)  node[below right]{$\sigma_3$};
	
	\coordinate (o) at (0,0);
	\coordinate (i) at (4,0);
	\coordinate (j) at (0,4);
	
	\coordinate (u) at (2.1,2.14);
	\coordinate (p) at (2.1, -0.5);
	\coordinate (v) at (3.4,0.5);
	
	% cercles
	\draw[thick] (4,0) arc [start angle=0, end angle =-140, radius = 4];
	\draw[thick] (4,0) arc [start angle=0, end angle =100, radius = 4];
	
	\draw[thick, dashed] (4,0) arc [start angle=0, end angle =100, radius = 4, y radius = 1];
	\draw[thick] (4,0) arc [start angle=0, end angle =-125, x radius = 4, y radius = 1];
	
	% droites
	\draw[-stealth,thick, blue] (0,0) -- (u) node[above right] {$\rho$};
	\draw[dashed, blue] (u) -- (p);
	\draw[dashed, blue] (0,0) -- (p);
	
	%arc au point
	\draw[gray, dashed] (3.1,2.5) arc [start angle=0, end angle =105, radius = 3.1, y radius = 0.5];
	\draw[gray] (3.1,2.5) arc [start angle=0, end angle =-120, x radius = 3.1, y radius = 0.5];
	
	\draw[violet, -stealth] (-0.4, -0.2) arc [start angle=-120, end angle =-40, x radius = 0.9, y radius = 0.225] node[midway, below] {\scriptsize$2\theta$};    
	\draw[violet, -stealth] (0,0.7) arc [start angle=90, end angle =20, x radius = 0.4, y radius = 0.5];
	\draw[violet] (0.25, 0.8) node{\scriptsize$2\chi$};
\end{tikzpicture}
	\caption{Sphère de Poincaré, blablabla}
	\label{fig:sphere2pointcare}
\end{wrapfigure}

\par
où les $\sigma_i$ sont les matrices de Pauli :
\begin{align*}
	\sigma_1 &= \begin{pmatrix} 0 & 1 \\ 1 &  0 \end{pmatrix}  &
	\sigma_2 &= \begin{pmatrix} 0 & -i \\  i &  0 \end{pmatrix}  &
	\sigma_3 &= \begin{pmatrix} 1 & 0 \\ 0 & -1 \end{pmatrix}
\end{align*}
\par

On remarque alors que les coefficients associés aux $\sigma_i$ sont des coordonnées sphériques lorsque les $\sigma_i$ sont associés à la base canonique de $\R^3$, conformément à la \Cref{fig:sphere2pointcare} ci-contre \thoughts{pas à jour}.
\\
Posé ainsi, l'intégrale dans la phase dynamique ne s'interprète pas. Vraiment pas, \href{https://en.wikipedia.org/wiki/Spherical_coordinate_system#Integration_and_differentiation_in_spherical_coordinates}{Wikipedia} liste une tonne de formule mais aucune ne colle avec $\theta'\sin2\chi$ parce que les angles choisis ne le permette pas. 
Notre convention est que $2\chi$ correspond à l'angle formé par $\sigma_2$ et $\rho$, là où sur Wikipedia, il est définie par l'angle entre la projection de $\rho$ sur $\Vec\{\sigma_{1,2}\}$ et $\rho$. Pour le dire rapidement, ce la change tout les $\cos 2\chi$ en $\sin 2\chi$ et inversement (modulo des signes).
\\
Pour pouvoir interpréter l'intégrale, il faut choisir les mêmes conventions qu'eux, auquel cas, il faut poser :
\[\rho = \sin2\theta\sin2\chi \sigma_1 + \cos 2\chi \sigma_3 + \cos2\theta \sin2\chi \sigma_3\]
\\
et notre vecteur apparaît comme la dérivée :
\[\hat{\bf{\chi}} = \frac{\partial}{\partial\chi} \rho = \sin2\theta\cos2\chi \sigma_1 - \sin 2\chi \sigma_3 + \cos2\theta \cos2\chi \sigma_3\]
\\
Dans ce cas, $\theta'\sin2\chi$ s'interprète (à un facteur près) comme la norme de la composante horizontale du vecteur vitesse de $\rho$, puisque :
\[\partial_t\rho = 2 \chi'\hat{\bf{\chi}} + 2 \theta' \sin 2\chi \hat{\bf{\theta}}\]
\\
Ainsi, mais c'est vraiment pas beau, l'intégrale de $\phased$ s'interprète comme la longueur du déplacement horizontale-ish de la projection $\rho$ de $\x$.
\\
Si on veut le voir en terme d'élément d'air $dS$, alors on a :
\[dS = \sin2\chi d\theta d\chi\]
\\
Mais comme $\theta$ et $\chi$ dépendent d'un paramètre commun $t$, ça serait très impropre de parler de calcul d'air sachant que $\theta'(t)\sin 2\chi(t)$ parcours un chemin vraiment. D'autant plus que là pas d'intégrale sur $\theta$... même si peut-être qu'on pourrait la faire apparaître avec un changement de variable approprié (surement que non en plus parce que le passage de $t$ à $(\theta,\chi)$ à par une différentielle inversible donc au cul le det de la Jacobienne :/ )

\subsection{Généralisation en plus haute dimension} \label{subsec:gene_AM-FM-PM}

\begin{itemize}
	
	\item Différentes écritures du bivarié pour différentes généralisation :
	
	\item Les quaterions on passe vites parce que ca se généralise très mal, Lefevre a a parlée, ca mène aux algèbres Clifford : trop de contrainte sur les dimensions des signaux
	
	\item En terme d'expo de matrice ? Lefevre \cite[sec. I.3]{lefevre_polarization_2021} l'a fait en trivarié mais au delà, y'a plus vraiment de choix remarquable de base pour $\mathfrak{u}(n)$
	
	\item En augmentant la taille de la matrice de rotation ? Lilly \cite{lilly_modulated_2011} l'a fait en trivarié et mais là encore, en terme de généralisation c'est pas si dingue parce que la matrice de rotation est pas calculable.
	
\end{itemize}
	


\subsection{Sur la nécessité d'avoir une approche géométrique}\label{subsec:aller_plus_loin}
\begin{itemize}
	
	\item Dans tout ça, on ratte le plus important : La phase géo est invariante par transfo de jauge, donc il faut faut faire apparaître $\PC{n-1}$ dans la décomposition.
	
	\item et en fait, c'est le cas en bivarié car $\PC{1}\cong \S{2}$ !
	
	\item $\PC{n-1}$ oui mais il faut pas non plus regarder que la projection parce qu'on perd toute les phases dans ce cas.
	
	\item Le bon compromis c'est les variétés fibrées : on est dans $\PC{n-1}$ mais on garde les phases dans les fibres.
	
	\item D'autant plus que ça à déjà était fait en physique et c'est vraiment concluant... (transition vers la grande partie suivante.)
	
\end{itemize}





%%%% ANNEXES %%%%



\begin{annexe}

\section{Annexes}

\subsection{Compléments sur l'analyse temps-fréquence}\label{ann:complement_t-f}

\subsubsection{Un mot sur la notion de fréquence instantanée \thoughts{(nécessaire?)}}\label{ann:freq_instant}

\subsubsection{Formalisme dernière la transformée en SA}\label{ann:transfo_SA}

\subsubsection{Lien avec le théorème de Bedrosian}\label{ann:bedrosian}

\subsection{Compléments sur les signaux AM-FM-PM}\label{ann:AM-FM-PM}

\subsubsection{Construction détaillée des signaux AM-FM-PM \thoughts{(nécessaire?)}}\label{ann:construc_AM-FM-PM}

\begin{itemize}
	
	\item Signal polarisé classique ($\theta,\chi$ constants)
	
	\item Transformé en SA avec les hypothèse de Bedrosian 
	
	\item Définition général en faisant dépendre $\theta / \chi$ du temps
	
	\item Paramètre de Polar sur la sphère de Poincaré
	
\end{itemize}


\subsubsection{Démostrations}\label{ann:demo_phases_2var}

\begin{demo}[de la formule \eqref{eq:phased_2var}, \cref{prop:phases_2var}]
	Par souci de lisibilité, on note $\mathcal{U} = R_{\theta} \begin{pmatrix} \cos\chi \\ -i\sin\chi \end{pmatrix}$ de sorte que la dérivée de $\x$ s'écrive :
	\begin{align*}
		\dot{\x} 
		&= a'e^{i\varphi}\mathcal{U} + ia\varphi'e^{i\varphi} \mathcal{U} + ae^{i\varphi}\theta'\begin{pmatrix} -\sin\theta \cos\chi + i\cos\theta \sin\chi \\ \cos\theta \cos\chi + i\sin\theta \sin\chi \end{pmatrix} + ae^{i\varphi}\chi'\begin{pmatrix} -\cos\theta \sin\chi + i\sin\theta \cos\chi \\ -\sin\theta \sin\chi - i\cos\theta \cos\chi \end{pmatrix} \\
		&= a'e^{i\varphi}\mathcal{U} + ia\varphi'e^{i\varphi} \mathcal{U} + ae^{i\varphi}\theta'\begin{pmatrix} 0 & -1 \\ 1 & 0 \end{pmatrix}\mathcal{U} + ae^{i\varphi}\chi'\begin{pmatrix} 0 & i \\ -i & 0 \end{pmatrix}\congu{\mathcal{U}}
	\end{align*}
	\\
	Le produit hermitien $\langle \x, \dot{\x}\rangle$ s'écrit alors :
	\begin{align*}
		\langle \dot{\x}, \x \rangle 
		&= \left\langle a'e^{i\varphi}\mathcal{U} + ia\varphi'e^{i\varphi} \mathcal{U} + ae^{i\varphi}\theta'\begin{pmatrix} 0 & -1 \\ 1 & 0 \end{pmatrix}\mathcal{U} + ae^{i\varphi}\chi'\begin{pmatrix} 0 & i \\ -i & 0 \end{pmatrix}\congu{\mathcal{U}}, ae^{i\varphi}\mathcal{U} \right\rangle \\
		&= \left\langle a'\mathcal{U} + ia\varphi' \mathcal{U} + a\theta'\begin{pmatrix} 0 & -1 \\ 1 & 0 \end{pmatrix}\mathcal{U} + a\chi'\begin{pmatrix} 0 & i \\ -i & 0 \end{pmatrix}\congu{\mathcal{U}}, a\mathcal{U} \right\rangle \\
		&= aa' \big\langle \mathcal{U}, \mathcal{U}\big\rangle  + ia^2\varphi' \big\langle \mathcal{U}, \mathcal{U}\big\rangle  + a^2\theta'\left\langle \begin{pmatrix} 0 & -1 \\ 1 & 0 \end{pmatrix}\mathcal{U}, \mathcal{U} \right\rangle - ia^2\chi'\left\langle \begin{pmatrix} 0 & -1 \\ 1 & 0 \end{pmatrix}\congu{\mathcal{U}}, \mathcal{U} \right\rangle
	\end{align*}
	où les deux derniers produits hermitiens donnent :
	\begin{align*}
		\left\langle \begin{pmatrix} 0 & -1 \\ 1 & 0 \end{pmatrix}\mathcal{U}, \mathcal{U} \right\rangle &= -\mathcal{U}_2\congu{\mathcal{U}_1} + \mathcal{U}_1\congu{\mathcal{U}_2} \\
		&= 2i\Im m\big(\mathcal{U}_1 \congu{\mathcal{U}_2}\big) \\
		&= 2i\Im m\Big(\big( \cos\theta \cos\chi + i \sin\theta \sin\chi \big) \big( \sin\theta \cos\chi + i \cos\theta \sin\chi \big)\Big) \\
		&= 2i\big(\cos^2\theta \cos\chi \sin\chi + \sin^2\theta \sin\chi \cos\chi \big) \\
		&= 4i \cos\chi \sin\chi \\
		&= 2i\sin2\chi 
		\\ \\
		\left\langle \begin{pmatrix} 0 & -1 \\ 1 & 0 \end{pmatrix}\congu{\mathcal{U}}, \mathcal{U} \right\rangle &= -\congu{\mathcal{U}_2\mathcal{U}_1} + \congu{\mathcal{U}_1\mathcal{U}_2} = 0
	\end{align*}
	\\
	D'où, sachant que $\ \|\x\|^2=a^2\ $ et $\ \|\mathcal{U}\|=1$, la formule :
	\begin{align*}
		\frac{\Im m\big\langle \dot{\x}, \x\big\rangle}{\|\x\|^2} &= \frac{1}{a^2}\Im m\Big(aa' \big\langle \mathcal{U}, \mathcal{U}\big\rangle  + ia^2\varphi' \big\langle \mathcal{U}, \mathcal{U}\big\rangle + ia^2\theta' \sin2\chi \Big) \\
		&= \frac{1}{a^2} \Big( a^2\varphi' \|\mathcal{U}\|^2 + a^2\theta' \sin2\chi \Big) \\
		&= \varphi' + \theta' \sin2\chi
	\end{align*}
\end{demo}

\begin{demo}[de la formule \eqref{eq:phaset_2var}, \cref{prop:phases_2var}]
	Pour la phase totale, on note cette fois $\mathcal{V} = \begin{pmatrix} \cos\chi \\ -i\sin\chi \end{pmatrix}$ et on a :
	\begin{align*}
		\big\langle \x(t), \x(t_0)\big\rangle &= \Big\langle a(t)e^{i\varphi(t)}R_{\theta(t)}\mathcal{V}(t), a(t_0)e^{i\varphi(t_0)}R_{\theta(t_0)}\mathcal{V}(t_0) \Big\rangle \\
		&= a(t)e^{i\varphi(t)} a(t_0)e^{-i\varphi(t_0)} \Big\langle R_{\theta(t)}\mathcal{V}(t), R_{\theta(t_0)}\mathcal{V}(t_0) \Big\rangle \\
		%&= a(t_0)a(t)e^{i(\varphi(t_0)-\varphi(t))}\Big\langle \mathcal{V}(t_0), R_{\theta(t_0)}^{-1}R_{\theta(t)}\mathcal{V}(t) \Big\rangle \\
		&= a(t_0)a(t)e^{i(\varphi(t)-\varphi(t_0))}\Big\langle R_{\theta(t)- \theta(t_0)}\mathcal{V}(t), \mathcal{V}(t_0) \Big\rangle
	\end{align*}
	Pour alléger les notations, on note $\ \Delta y =y(t)-y(t_0)$, $\ y_1=y(t_0)\ $ et $\ y_2=(t)\ $ pour $\ y=\varphi,\theta,\chi$. Le produit hermitien à droite s'écrit alors :
	\begin{align*}
		\Big\langle R_{\Delta\theta}\mathcal{V}(t), \mathcal{V}(t_0) \Big\rangle &=   \Big(\cos\Delta\theta \cos\chi_2 + i\sin\Delta\theta \sin\chi_2 \qquad \sin\Delta\theta \cos\chi_2 - i\cos\Delta\theta \sin\chi_2\Big) \begin{pmatrix} \cos\chi_1 \\ i\sin\chi_1 \end{pmatrix} \\
		&= \cos\chi_1\Big(\cos\Delta\theta \cos\chi_2 + i\sin\Delta\theta \sin\chi_2\Big) + i\sin\chi_1\Big(\sin\Delta\theta \cos\chi_2 - i\cos\Delta\theta \sin\chi_2\Big) \\
		&= \cos\Delta\theta \Big(\cos\chi_1 \cos\chi_2 + \sin\chi_1 \sin\chi_2\Big) + i\sin\Delta\theta \Big( \cos\chi_1 \sin\chi_2 + \sin\chi_1\cos\chi_2\Big) \\
		&= \cos\Delta\theta \cos\Delta\chi + i\sin\Delta\theta \sin(\chi_1+\chi_2)
	\end{align*}
	\\
	D'où la phase totale :
	\begin{align*}
		\phaset(\x) = \arg\big\langle \x(t), \x(t_0)\big\rangle &= \arg\left( a(t_0)a(t)e^{i(\varphi(t) - \varphi(t_0))}\Big( \cos\Delta\theta \cos\Delta\chi + i\sin\Delta\theta \sin(\chi_1+\chi_2) \Big) \right) \\
		&= \varphi(t) - \varphi(t_0) + \arg\Big( \cos\Delta\theta \cos\Delta\chi + i\sin\Delta\theta \sin(\chi_1+\chi_2) \Big)
	\end{align*}
	et l'argument restant s'écrit comme une arctangente, donnant :
	\begin{align*}
		\phaset(\x) &= \varphi(t)-\varphi(t_0) + \arctan\frac{\sin\Delta\theta \sin(\chi_1+\chi_2)}{\cos\Delta\theta \cos\Delta\chi} \\
		&= \varphi(t)-\varphi(t_0) + \arctan \left( \tan\Delta\theta\frac{\sin(\chi_1+\chi_2)}{\cos\Delta\chi} \right) \\
		&= \cdots
	\end{align*}
\end{demo}

\end{annexe}